% Round 16. New actual-base-change and exact-support strong-G statements.
% Activated after independent trace-branch review; no converse G => obstruction zero.
\section{Projective-bundle base change and exact strong Green stability}
\label{r16pg:section}

The coordinate base-change theorem OUT16.16 extends to an arbitrary
algebraic testing map from a smooth variety. The finite-flat identity
H16.1, rather than a base-change rule for arbitrary currents, supplies
the additional descent step. This yields a statement about an arbitrary
strong Green form; its hypothesis need not be certified by a vanishing
actual Bott--Chern obstruction.
The mathematical source passed independent trace-branch review at
SHA256 \texttt{2a917991d0511756587e8a834feeb586fd509346bcc5a9230b259101dd9e219e};
the audit is
\path{round16/higher_trace/projective_strong_green_base_change_independent_audit.md}.
Only this activation notice differs from that reviewed version.

\begin{theorem}[Actual projective-bundle integration under a smooth testing map]
\label{r16pg:base-change}
Let $X/K$ be smooth geometrically integral, let $N$ be an algebraic
vector bundle of rank $m+1$, and let $\pi:E=\mathbf P_X(N)\to X$.
Let $a:M\to X$ be any algebraic morphism with $M$ smooth integral,
and form the Cartesian square with $E_M$, $\pi_M$ and $a_E$.
For every ordinary open $U\subset X^{\rm an}$ and every
$\alpha\in B^{p,q}(E_U)$, $p,q\geq0$, one has
\[
 a^*(I_\pi\alpha)=I_{\pi_M}(a_E^*\alpha)
 \quad\text{on }(a^{\rm an})^{-1}U.
 \tag{PG16.1}\label{r16pg:bc}
\]
Both sides are ordinary actual forms; negative output degrees mean
zero. The map $a$ need not be smooth, flat, dominant or an immersion.
The construction retains the original $K$-weights and full degrees.
\end{theorem}

\begin{proof}
First take $E=X\times\mathbf P^1$. Fix $y\in M^{\rm an}$ and
choose one original AD construction for $\alpha$ around $a(y)$.
Keep its entire finite etale splitting cover
$q:U'\to U_0$, of full degree $e$, all its root functions and
difference units, its neutral horizontal cutoff $\rho$, and its
supplied split presentation
\[
 \rho\alpha'=t_{(b,w)}^\star\omega,
 \qquad
 \beta'=t_b^\star\operatorname{pr}_*(\Gamma_s\wedge\omega).
 \tag{PG16.2}\label{r16pg:supplied}
\]
The ordinary opens here are restrictions of the algebraic finite
etale cover in AD. If that cover is disconnected, retain each of
its integral components and sum their traces. Its total degree
is $e$. Its output downstairs is $e^{-1}\operatorname{Tr}_q\beta'$.

Pull this entire data to $M$. Around $y$ choose a compact working
neighborhood inside $a^{-1}U_0$; there is no claim that the full
inverse image of the original compact neighborhood is compact.
The pulled cover is finite etale, hence its components are smooth
and already normal, and its total degree is still $e$.
Every difference unit remains a unit, every strict numerical
domain margin is preserved, and the pulled roots remain distinct.
The actual identity in \eqref{r16pg:supplied} pulls back literally.
The same cutoff removes tubes on which the pulled input is a
pullback from the new base: this follows by pulling the original
retraction-germ identity. Its neutral-excision proof therefore
works with every actual target factor on that new base, by OUT16.2.

Apply the supplied-presentation UG theorem on each component of
the pulled cover. It computes the output using exactly the same
Euclidean coefficient $\operatorname{pr}_*(\Gamma_s\wedge\omega)$,
now pulled through the base units $a^*b$. No full-rank condition
on this initially supplied unit tuple is required. Append all
new target test and localization moments before the refined-cycle
calculation. The refined base cycle lies in the same universal
root-difference fan times the extra base-coordinate space; UG15.4
applies to every such cycle with its native weight. The source
has relative weight one by FT15.7. Thus the identical coefficient
has all the prescribed smooth-test pairings on the new base, and
OUT16.1--OUT16.6 give the full actual-test and target-factor
identities there. This is an actual representative construction,
not restriction of a pushforward current to the image of $a$.
The degree-zero-output case uses the same ambient formula and
is covered by OUT16.3.

It remains to descend. Apply H16.1 to each integral component of
the original finite etale cover and the smooth testing map $a$.
Every pulled component has generic length one; the sum of all
its native degrees is the original degree. Consequently
\[
 a^*\bigl(e^{-1}\operatorname{Tr}_q\beta'\bigr)
   =e^{-1}\sum_j\operatorname{Tr}_{q_j}(a_j^*\beta').
 \tag{PG16.3}\label{r16pg:descent}
\]
The right side is exactly the local AD/UG formula for the new
projective-line output. This proves \eqref{r16pg:bc} locally for
the line, hence everywhere by ordinary actual uniqueness.

For $X\times\mathbf P^m$, use the same finite symmetric map
$(\mathbf P^1)^m\to\mathbf P^m$ of full degree $m!$ and iterate
the line construction OUT16.12. At each stage the testing map
is $a$ times the identity on the remaining projective factors.
Its target is geometrically integral and its source is smooth
integral, so the proved line identity applies. Pullback commutes
with the symmetric map as a genuine algebraic Cartesian diagram;
the normalizing degree is unchanged. The rank-one case is the
identity. Algebraic local trivializations of $N$, restriction
compatibility, and actual uniqueness give the asserted bundle
formula. No step selects one sheet of a finite cover.
\end{proof}

\begin{lemma}[Restriction of the actual output to a cycle]
\label{r16pg:cycle}
Assume $X$ and $E$ are projective. For a specified rational
codimension-$r$ cycle $Z$ on $X$, put $Z_E=\pi^*Z$ (flat
algebraic pullback). Then for every actual form
$\eta\in B^{n+m-r,n+m-r}(E^{\rm an})$, where $n=\dim X$,
\[
                  \int_Z I_\pi\eta=\int_{Z_E}\eta.
 \tag{PG16.4}\label{r16pg:cycle-eq}
\]
The formula uses the specified prime coefficients, including
native constant-field degrees. It does not posit a projective-
bundle actual operator on a singular base.
\end{lemma}

\begin{proof}
It suffices to take $Z$ integral. Let $Z^\circ$ be its dense
smooth algebraic locus. The restriction of $I_\pi\eta$ to $Z$
is an actual top form; the restriction of $\eta$ to
$E_Z=E\times_X Z$ is also top degree. By top-degree support
avoidance both forms vanish on ordinary neighborhoods of the
respective proper algebraic closed subsets $Z\setminus Z^\circ$
and $E_Z\setminus E_{Z^\circ}$. Their supports are therefore
compact subsets of the displayed smooth loci, since $Z$ and
$E_Z$ are proper. The map $Z^\circ\to X$ is a testing map
permitted by Theorem~\ref{r16pg:base-change}; it identifies the
first top form with the actual projective-bundle output of the
second on $Z^\circ$. Use OUT16.11 there, inserting a compact
ordinary smooth scalar equal to one near the proper projection
of the second support and the first support. This gives
\eqref{r16pg:cycle-eq}. Smooth geometrically integral projective
fibers make $E_Z$ integral, with generic flat multiplicity one.
Linearity proves the statement for rational cycles. No removal
changes a top integral or a native cycle coefficient.
\end{proof}

\begin{theorem}[Exact strong G is stable in both directions under a projective bundle]
\label{r16pg:strong-G}
Let $X/K$ be smooth projective geometrically integral of dimension
$n$, let $\pi:E=\mathbf P_X(N)\to X$ have relative dimension $m$,
and let $Z$ be a specified rational cycle of codimension
$1\leq r\leq n$. Put $Z_E=\pi^*Z$. Then
\[
                         G_X(Z)\quad\Longleftrightarrow\quad G_E(Z_E).
 \tag{PG16.5}\label{r16pg:G-equivalence}
\]
Here both statements mean ordinary actual Green forms on the exact
specified support complement, with zero curvature against every
compact actual test over the original $K$. No obstruction-vanishing
hypothesis is imposed.

More precisely, if $\gamma_Z$ is any such form downstairs, then
$\pi^*\gamma_Z$ is one upstairs. Conversely, choose any PWS metric
on $\mathcal O_E(1)$, put $\xi=c_1(\overline{\mathcal O_E(1)})$,
and, for any strong form $\gamma_{Z_E}$ upstairs, put
\[
      \gamma_Z=I_{\pi|_{X\setminus|Z|}}
                         (\xi^m\wedge\gamma_{Z_E}).
 \tag{PG16.6}\label{r16pg:descent-G}
\]
This is a strong form downstairs. For $m=0$ both constructions are
the identity, with $\xi^0=1$.
\end{theorem}

\begin{proof}
The flat inverse image of an integral prime cycle is integral,
and inverse images of distinct primes are distinct. Hence
$|Z_E|=\pi^{-1}|Z|$ even for signed rational coefficients written
after cancellation. This proves that genuine actual pullback and
the displayed bundle operator have exactly the required domains.

For ascent take an arbitrary compact actual test
$\eta\in B_c^{n+m-r,n+m-r}(E^{\rm an})$. The support of $D\eta$
avoids $|Z_E|$ by degree. Its proper image is a compact subset of
$X\setminus|Z|$, and the support property of the actual operator
places $I_\pi D\eta$ there. The projection formula, full actual
adjunction and differential commutation give
\[
 \begin{split}
 \int_{E\setminus|Z_E|}\pi^*\gamma_Z\wedge D\eta
   &=\int_{X\setminus|Z|}\gamma_Z\wedge I_\pi D\eta\\
   &=\int_{X\setminus|Z|}\gamma_Z\wedge D(I_\pi\eta)\\
   &=-\int_Z I_\pi\eta
    =-\int_{Z_E}\eta.
 \end{split}
 \tag{PG16.7}\label{r16pg:ascent-pairing}
\]
For the first equality insert an ordinary smooth compact scalar
equal to one near the projection of $\operatorname{supp}D\eta$;
this makes the target factor a permitted compact actual test.
The strong identity downstairs applies to the global compact
actual form $I_\pi\eta$. The last equality is
Lemma~\ref{r16pg:cycle}. Thus the assertion is on the full actual
test space, not just tests pulled back from $X$.

For descent, the frozen Chern-degree lemma PC6 applied to the
projective map $\pi$ and $\xi^m$ gives
\[
                             I_\pi(\xi^m)=1.
 \tag{PG16.8}\label{r16pg:unit}
\]
Indeed, its pairing with every compact actual top test is the
algebraic pushforward
$\pi_*(c_1(\mathcal O_E(1))^m\cap[E])=[X]$, and actual-current
injectivity identifies the scalar output with one. Both $d'\xi$
and $d''\xi$ vanish. For an arbitrary
$\eta\in B_c^{n-r,n-r}(X^{\rm an})$, bundle adjunction and the
upstairs strong identity therefore give
\[
 \begin{split}
 \int_{X\setminus|Z|}\gamma_Z\wedge D\eta
   &=\int_{E\setminus|Z_E|}\gamma_{Z_E}
                                  \wedge D(\xi^m\wedge\pi^*\eta)\\
   &=-\int_{Z_E}\xi^m\wedge\pi^*\eta\\
   &=-\int_Z I_\pi(\xi^m\wedge\pi^*\eta)
    =-\int_Z\eta.
 \end{split}
 \tag{PG16.9}\label{r16pg:descent-pairing}
\]
The source test in the first line has degree
$(n+m-r,n+m-r)$ and is compact by projectivity. Its derivative
has support outside the polar set. The last two equalities use
Lemma~\ref{r16pg:cycle}, the all-actual projection formula, and
\eqref{r16pg:unit}. This also verifies every degree and support
requirement for the displayed actual representative.
\end{proof}

\begin{corollary}[Preservation of the finite-intersection P property by flat pullback]
\label{r16pg:P}
In Theorem~\ref{r16pg:strong-G}, suppose the chosen $\gamma_Z$
satisfies the finite-intersection identity H16.7 for every normal
integral quasiprojective $r$-dimensional testing map to $X$.
Then $\pi^*\gamma_Z$ satisfies the same property for $Z_E$ on $E$.
The same assertion holds for the stronger proper-intersection
identity HP16.2 whenever that identity is supplied for $\gamma_Z$.
\end{corollary}

\begin{proof}
For a testing map $a:M\to E$, compose with $\pi$. Its inverse
image of $|Z|$ is exactly $a^{-1}|Z_E|$. Pullback of the Green
form is literally $(\pi a)^*\gamma_Z$. Refined Gysin composition
and smooth flat pullback identify
$I_a(\pi^*Z)=I_{\pi a}(Z)$, with the same scheme lengths. Apply
the supplied identity to the composite. This proof does not
deduce a P property from a bare strong current equation.
\end{proof}

\paragraph{Dependencies and limitations.}
The only new base-change interface used is H16.1 for finite flat
actual trace. The supplied ambient construction is Round 15
\path{root/universal_graphic_tree_kernel.tex}, UG15.4 and UG15.8,
with \path{actual_local/finite_ambient_delta_atlas.tex}, AD15.3.
The full actual outputs, arbitrary actual target-factor projection,
zero-output convention, and full test comparison are OUT16.1--16
in Round 16 \path{higher_trace/all_actual_projective_bundle_output.tex}.
PC6 is the frozen Round 11 supported-complex Chern-degree theorem.
No general proper actual pushforward, arbitrary current pullback,
or implication from G to obstruction vanishing is used. These
projective-bundle statements alone do not resolve the remaining
primitive classes on a general variety.
